\section{Two horizons: de Sitter QNMs and continuum response}
\label{sec:bh-desitter}
% BEGIN CURATED SUBJECT INDEX
\index{quasinormal mode}
\index{horizon boundary condition}
\index{de Sitter spacetime}
\index{continuum response}
% END CURATED SUBJECT INDEX

\subsection{Two-horizon scattering problem}
\label{subsec:sds-setup}
% BEGIN CURATED SUBJECT INDEX
\index{boundary condition!incoming and outgoing}
\index{complex scaling}
\index{quasinormal mode}
\index{tortoise coordinate}
\index{horizon boundary condition}
\index{Schwarzschild black hole}
\index{de Sitter spacetime}
% END CURATED SUBJECT INDEX

For a four-dimensional Schwarzschild--de Sitter black hole,
\begin{equation}
  f(r)=1-\frac{2M}{r}-\frac{\Lambda r^2}{3},
  \qquad \Lambda>0 .
  \label{eq:sds-f}
\end{equation}
In the static region, $f(r)$ has a black-hole horizon $r_h$ and a cosmological
horizon $r_c>r_h$.  The tortoise coordinate maps
\begin{equation}
  r\to r_h^+\Longrightarrow r_*\to-\infty,
  \qquad
  r\to r_c^-\Longrightarrow r_*\to+\infty .
  \label{eq:sds-tortoise-limits}
\end{equation}
The QNM boundary conditions are
\begin{align}
  \psi(r_*)&\sim\rme^{-\rmi\omega r_*},
  &&r_*\to-\infty,
  \notag\\
  \psi(r_*)&\sim\rme^{+\rmi\omega r_*},
  &&r_*\to+\infty .
  \label{eq:sds-boundary}
\end{align}
They have the same two-ended outgoing form as Eq.~\eqref{eq:qnm-boundary-flat},
with the right end now representing the cosmological horizon.

For scalar, electromagnetic, and axial gravitational sectors, one may again
use
\begin{equation}
  V_{sl}(r)
  =f(r)\left[
  \frac{l(l+1)}{r^2}
  +\frac{f'(r)}{r}\,\delta_{s0}
  -\frac{6M}{r^3}\,\delta_{s2}
  \right] ,
  \label{eq:sds-potentials}
\end{equation}
with the sector labels interpreted appropriately.  The potential decays
exponentially at both horizons for a nonextremal static patch.  This
short-range behavior makes global complex scaling in $r_*$ particularly
natural, subject to the analytic branches of $r(r_*)$.

\subsection{Pole motion with the cosmological constant}
\label{subsec:sds-pole-motion}
% BEGIN CURATED SUBJECT INDEX
\index{boundary condition!incoming and outgoing}
\index{quasinormal mode}
\index{horizon boundary condition}
% END CURATED SUBJECT INDEX

Changing $\Lambda$ changes the horizon separation, surface gravities, barrier
height, and barrier width.  Quasinormal poles therefore move in the complex
$\omega$ plane.  A complex-scaling calculation follows these pole trajectories
without changing the outgoing boundary implementation.  The same diagonalized
spectrum also contains the rotated continuum, allowing pole motion to be
distinguished from redistribution of background response
\cite{Ogawa:2026dSCSM}.

Near the Nariai limit, the two horizons approach one another and an appropriate
rescaled potential approaches a solvable barrier.  This gives an analytic
benchmark related to the P\"oschl--Teller family.  Away from that limit, a
numerical comparison should keep fixed a dimensionless convention, such as
$M\omega$ or a frequency scaled by a surface gravity; otherwise apparent
motion partly reflects a changing unit.

\subsection{Continuum level density for black-hole response}
\label{subsec:bh-cld}
% BEGIN CURATED SUBJECT INDEX
\index{scattering matrix}
\index{continuum level density}
\index{horizon boundary condition}
\index{continuum response}
% END CURATED SUBJECT INDEX

Choose a reference operator $H_0$ that has the same horizon asymptotics as the
black-hole master operator but no scattering barrier.  The CLD is
\begin{equation}
  \Delta\rho(E)
  =-\frac{1}{\pi}\im\Tr\left[
  (E-H_\theta)^{-1}
  -(E-H_{0,\theta})^{-1}
  \right],
  \qquad E=\omega^2 .
  \label{eq:bh-cld}
\end{equation}
A stable peak can be decomposed locally into a pole contribution and a smooth
background, but the total CLD is the more invariant finite quantity.  Its
integral measures the total scattering phase relative to the reference.  It
therefore probes changes in continuum response even when no isolated pole
crosses the observation window.

For a two-ended one-dimensional potential, the scattering matrix contains
reflection and transmission amplitudes.  The determinant phase fixes the CLD
through Eq.~\eqref{eq:birman-krein-differential}, while the greybody factor is
the transmission probability
\begin{equation}
  \mathcal G_l(\omega)=|T_l(\omega)|^2 .
  \label{eq:greybody-factor}
\end{equation}
Knowing $\det S$ does not in general determine $|T|^2$.  Additional symmetry
or channel-resolved information is required.  The CLD and greybody factor are
therefore complementary rather than interchangeable observables
\cite{Harmark:2007jy,Ogawa:2026dSCSM}.

\subsection{Coupled channels and string-inspired backgrounds}
\label{subsec:bh-coupled-channels}
% BEGIN CURATED SUBJECT INDEX
\index{spectrum!point}
\index{scattering matrix}
\index{boundary condition!incoming and outgoing}
\index{Riemann sheet}
\index{resonance!residue}
\index{complex scaling}
\index{coupled-channel equation}
\index{horizon boundary condition}
% END CURATED SUBJECT INDEX

For $N_c$ coupled perturbation channels, write
\begin{equation}
  \left[-\identity_{N_c}\frac{\rmd^2}{\rmd r_*^2}
  +\bm V(r)\right]\bm\psi(r_*)
  =\omega^2\bm\psi(r_*) .
  \label{eq:coupled-bh-equation}
\end{equation}
Complex scaling acts on the common coordinate, while $\bm V$ mixes channels.
The asymptotic channel basis must be chosen so that incoming and outgoing
directions are defined at each horizon.  If different channels have different
thresholds or asymptotic wave numbers, the rotated spectrum contains several
rays and the Riemann sheet must be labeled channel by channel.

String-inspired metric--dilaton systems provide a natural application
\cite{Bian:2026crb}.  Numerically, the main new issue is matrix-block imbalance:
one channel or basis sector can acquire exponentially larger matrix elements
under deformation.  Similarity balancing can improve conditioning without
changing eigenvalues, but left--right residues must be transformed
consistently.

\subsection{Higher dimensions}
\label{subsec:bh-higher-d}
% BEGIN CURATED SUBJECT INDEX
\index{horizon boundary condition}
% END CURATED SUBJECT INDEX

Tensor, vector, and scalar gravitational perturbations of higher-dimensional
static black holes also reduce to Schr\"odinger-type master equations in many
sectors
\cite{Kodama:2003jz,Ishibashi:2003ap,Ishibashi:2011ws}.  When the potential is
frequency independent and has two asymptotically free ends, the same
complex-scaling construction applies.  The angular spectrum and near-horizon
indices change, but the pole/continuum geometry does not.  Scalar-type coupled
systems and rotating backgrounds require additional work.
